\section{Proof of reflection for Boolean connectives}
\label{sec:connectives}

We give the intuitionistic argument in detail because this is where an
unmarked appeal to classical logic often enters a packet.  Throughout, suppose
\(a,b\) are Boolean terms and
\[
  a=1\leftrightarrow A,
  \qquad
  b=1\leftrightarrow B,
\]
with the corresponding zero/falsity equivalences supplied by
\cref{lem:atomic-zero} or the induction hypothesis.

\subsection{Negation}

Let \(\mathsf{notb}(a)=1\dotminus a\), where \(\dotminus\) is truncated
subtraction.  On the two Boolean cases,
\[
  \mathsf{notb}(0)=1,
  \qquad
  \mathsf{notb}(1)=0.
\]
Booleanity of \(\mathsf{notb}(a)\) follows by eliminating the proof
\(a=0\vee a=1\) and rewriting in each branch.

If \(\mathsf{notb}(a)=1\), the branch \(a=1\) is arithmetically impossible, so
\(a=0\), hence \(\neg A\).  Conversely, from \(\neg A\), falsity reflection
gives \(a=0\), and therefore \(\mathsf{notb}(a)=1\).  This proves truth
reflection for \(\neg A\); the zero half follows by the general Booleanity
lemma.

\subsection{Conjunction}

Define \(\mathsf{andb}(a,b)=a\cdot b\).  Four object-level cases on the two
Boolean disjunctions show that the product is Boolean.  If the product equals
one, neither factor can be zero; eliminating their Boolean cases gives
\(a=b=1\), hence \(A\wedge B\).  Conversely \(A\wedge B\) yields
\(a=b=1\) and the product equation.

This proof can be shortened with existing lemmas about products and zero, but
the four-case proof is useful as a generic implementation pattern: it invokes
only the reflected finite disjunctions and arithmetic rewriting.

\subsection{Disjunction}

Let \(\mathsf{orb}(a,b)=\mathsf{sg}(a+b)\), where \(\mathsf{sg}(0)=0\) and
\(\mathsf{sg}(Sx)=1\).  Again the four Boolean cases show that
\(\mathsf{orb}(a,b)=1\) exactly when at least one input is one.  Therefore
\[
  \mathsf{orb}(a,b)=1\leftrightarrow A\vee B.
\]
The converse direction is constructive: eliminate the given disjunction and
rewrite the corresponding bit to one.

An alternative implementation is
\(a+b\dotminus(a\cdot b)\).  The sign form is usually shorter because the
checked arithmetic library already exposes the zero/successor behavior of
\(\mathsf{sg}\).

\subsection{Implication}

Implication is the subtle connective.  Define
\[
  \mathsf{impb}(a,b)=\mathsf{orb}(\mathsf{notb}(a),b).
\]
We prove
\[
  \mathsf{impb}(a,b)=1\leftrightarrow(A\to B)
\]
without unrestricted excluded middle.

For the forward direction, assume the implication bit is one and assume
\(A\).  Truth reflection gives \(a=1\), hence
\(\mathsf{notb}(a)=0\).  The only way the disjunction bit can be one is then
\(b=1\), so \(B\).  Discharging the assumption gives \(A\to B\).

For the reverse direction, assume \(A\to B\) and eliminate the finite
Booleanity disjunction \(a=0\vee a=1\).  If \(a=0\), then
\(\mathsf{notb}(a)=1\), so the implication bit is one independently of \(b\).
If \(a=1\), truth reflection gives \(A\), the assumed implication gives
\(B\), and truth reflection for \(B\) gives \(b=1\).  Thus the implication bit
is again one.  The case split is justified by a characteristic bit already
inside iPRA; it does not assume that arbitrary formulas are decidable.

\subsection{Structural recursion and proof objects}

The host compiler recursively returns a record containing:
\[
  (\chi_A,\ \pi_A^{\mathrm{bool}},\ \pi_A^{1},\ \pi_A^{0}).
\]
At a connective it constructs a new fixed PR symbol by composition, then
combines the child derivations using ordinary natural-deduction constructors
and the checked arithmetic equations.  The recursion is over a finite Lean
datatype.  Its result is a syntax symbol and actual \code{Prf} terms, not a
semantic theorem about \(\Nat\).

\begin{proposition}[Connective fragment]
For every bounded-code expression formed from reflected atoms by
\(\neg,\wedge,\vee,\to\), the three conclusions (R0)--(R2) of
\cref{thm:reflection} hold.
\end{proposition}

\begin{proof}
Structural induction on the code.  The atom case is the stored package plus
\cref{lem:atomic-zero}; each connective case is the construction just given.
All generated terms are fixed compositions of primitive-recursive symbols.
\end{proof}
